\documentclass[11pt,a4paper]{article}

\usepackage[utf8]{inputenc}
\usepackage[T1]{fontenc}
\usepackage{geometry}
\geometry{left=1in,right=1in,top=1in,bottom=1in}
\usepackage{booktabs}
\usepackage{array}
\usepackage{enumitem}
\usepackage{hyperref}
\usepackage{amsmath,amssymb}
\usepackage{titlesec}

\hypersetup{
    colorlinks=true,
    linkcolor=blue,
    urlcolor=blue,
    pdftitle={Financial Fraud Detection Report}
}

\title{\textbf{Financial Fraud Detection with Graph Neural Networks and XGBoost}}
\author{}
\date{\today}

\begin{document}

\maketitle

\begin{abstract}
This report summarizes a fraud detection project on two graph-based datasets: the Elliptic Bitcoin transaction graph and the richer Elliptic++ transaction-address graph. We compare a tabular XGBoost baseline on Elliptic with several graph neural network models, including GCN variants and heterogeneous GNNs. The results show that the strongest model depends on how much relational structure is available: XGBoost performs best on the simpler Elliptic setting, while a multi-task directed heterogeneous GNN performs best on Elliptic++.
\end{abstract}

\section{Introduction}
Financial fraud is inherently relational. Fraudulent transactions often occur in groups, move through chains of accounts, or share address-level patterns that are difficult to capture with isolated tabular features. Graph neural networks are a natural fit because they can combine node attributes with neighborhood structure.

Because both datasets are highly imbalanced, the report emphasizes illicit-class F1, precision, recall, ROC-AUC, and PR-AUC when available. Accuracy is reported only where it helps interpret the class imbalance.

\section{Datasets}

\subsection{Elliptic Bitcoin Dataset}
The Elliptic dataset contains:
\begin{itemize}[noitemsep]
    \item 203,769 transactions
    \item 165 usable node features
    \item 234,355 directed edges, symmetrized to 468,710 undirected edges in the GCN experiment
    \item Label distribution: 157,205 unknown, 42,019 licit, 4,545 illicit
\end{itemize}

The split follows a temporal setup:
\begin{itemize}[noitemsep]
    \item Training: steps 1--34
    \item Testing: steps 35--49
\end{itemize}

This split is important because it reflects the real fraud setting where future behavior must be predicted from past behavior.

\subsection{Elliptic++ Dataset}
The Elliptic++ dataset extends the original transaction graph with address-level structure and wallet features:
\begin{itemize}[noitemsep]
    \item 203,769 transactions
    \item 182 baseline transaction features
    \item 234,355 transaction-to-transaction edges
    \item 837,124 transaction-to-address edges
    \item 2,868,964 address-to-address edges
    \item 1,268,260 wallet feature rows
    \item Label distribution: 157,205 unknown, 42,019 licit, 4,545 illicit
\end{itemize}

The Elliptic++ experiments also use a temporal split and evaluate address-level supervision, heterogeneous message passing, and multi-task learning.

\section{Models}

\subsection{XGBoost on Elliptic}
The tabular baseline treats fraud detection as a standard binary classification problem using transaction features only. It is a useful comparison point because it does not use graph structure.

\subsection{GCN on Elliptic}
The Elliptic GCN uses weighted loss to handle imbalance and predicts illicit transactions from the graph neighborhood plus node features.

\subsection{Graph Models on Elliptic++}
The Elliptic++ notebook compares:
\begin{itemize}[noitemsep]
    \item Baseline GCN
    \item Wallet-upgrade GCN
    \item Directed heterogeneous GNN
    \item Multi-task directed heterogeneous GNN
    \item Address pretraining followed by transaction fine-tuning
\end{itemize}

\section{Results}

\subsection{Elliptic}
\begin{table}[htbp]
\centering
\begin{tabular}{lccccc}
\toprule
Model & Illicit F1 & Precision & Recall & ROC-AUC & PR-AUC \\
\midrule
XGBoost & 0.7700 & 0.8100 & 0.7300 & 0.9230 & 0.7932 \\
GCN & 0.2410 & 0.1401 & 0.8596 & 0.8138 & N/A \\
\bottomrule
\end{tabular}
\caption{Elliptic test performance. Metrics are for the illicit class unless otherwise noted.}
\label{tab:elliptic_results}
\end{table}

The XGBoost baseline performs strongly on Elliptic, especially on illicit F1 and ROC-AUC. This suggests that the engineered transaction features already contain a lot of predictive signal.

The GCN achieves very high recall for illicit nodes, but precision is low. In practice, that means it finds many fraudulent transactions, but with many false positives.

\subsection{Elliptic++}
\begin{table}[htbp]
\centering
\begin{tabular}{lcccc}
\toprule
Model & Illicit F1 & Precision & Recall & ROC-AUC \\
\midrule
Baseline GCN & 0.2147 & 0.1325 & 0.5660 & 0.7772 \\
Wallet-upgrade GCN & 0.2245 & 0.1379 & 0.6038 & 0.7769 \\
Directed hetero GNN & 0.2250 & 0.1636 & 0.3601 & 0.7920 \\
Address pretrain + finetune & 0.2379 & 0.1502 & 0.5708 & 0.7885 \\
Multi-task directed hetero GNN & 0.2918 & 0.1875 & 0.6572 & 0.8560 \\
\bottomrule
\end{tabular}
\caption{Elliptic++ test performance. Metrics are for the illicit transaction class.}
\label{tab:ellipticpp_results}
\end{table}

The multi-task directed heterogeneous GNN is the best model on Elliptic++ by both illicit F1 and ROC-AUC. Relative to the baseline GCN, it improves by:
\begin{itemize}[noitemsep]
    \item +0.0771 in illicit F1
    \item +0.0788 in ROC-AUC
\end{itemize}

The address-pretraining strategy also improves over the baseline, but it does not match the multi-task model. The directed heterogeneous model improves ROC-AUC, but its recall is lower than the baseline GCN, so it is more conservative when flagging illicit nodes.

\section{Discussion}
The results point to three main conclusions:
\begin{enumerate}[noitemsep]
    \item Tabular models can still be very strong when the engineered feature space is rich. On Elliptic, XGBoost outperforms the GCN by a large margin.
    \item Graph structure becomes more valuable as the problem becomes more heterogeneous. On Elliptic++, the best model is the multi-task directed heterogeneous GNN.
    \item Threshold choice matters under severe class imbalance. Several GNN variants achieve reasonable ROC-AUC but still show low illicit precision, which means the ranking quality is decent but the decision threshold is not always optimal.
\end{enumerate}

\section{Limitations}
\begin{itemize}[noitemsep]
    \item Both datasets are highly imbalanced, so accuracy is not a reliable headline metric.
    \item Some models may improve with calibration or threshold tuning.
    \item Temporal drift may reduce performance when fraud behavior changes over time.
    \item The comparison focuses on the saved experimental outputs in the repository and excludes the logistic regression results and the \texttt{TestModels/} GCN results, as requested.
\end{itemize}

\section{Conclusion}
Across the two datasets, the best model depends on the amount of relational structure available.

\begin{itemize}[noitemsep]
    \item On the simpler Elliptic transaction graph, XGBoost is the strongest model.
    \item On the richer Elliptic++ graph, the multi-task directed heterogeneous GNN is the strongest model.
\end{itemize}

Overall, the project shows that graph learning becomes more useful as the dataset includes more node types, more edge types, and more auxiliary supervision.

\end{document}
